\documentclass{article}

\usepackage{mathtools}
\usepackage{amssymb}
\usepackage{unicode-math}
\setmonofont{DejaVu Sans Mono} % monospaced font with extensive unicode support for verbatim
\usepackage{stmaryrd}
\usepackage{booktabs}
\usepackage{tikz}
\usepackage{forest}
\usepackage{logicproof}
\usepackage{ebproof}
\usepackage{fitch}
\usepackage[capitalize,noabbrev]{cleveref}

%%%%%%% Some macros to help with the typesetting %%%%%%%%%%
\usepackage{logic-macros}
%%%%%% End of local macros %%%%%%%

\title{Logic Homework N}
\author{Name\\
  s1234567\\
  Studies%Informatica, I&E, Bioinformatica, Wiskunde, ...
}

\begin{document}

%% Use var symbols by default
\renewcommand*{\phi}{\varphi}
\renewcommand*{\epsilon}{\varepsilon}

\maketitle

This document is a template for the homework assignments.
It includes all packages necessary to produce truth tables, parse trees, and proofs.

\section*{Question 1}
\paragraph{a)}

Proofs written in Proof Rondo can be copy-pasted into a verbatim environment:
\begin{verbatim}
1 r and u → w   : premise
2 u and ~w      : premise
3 | r
4 | u           : and e1 2
5 | r and u     : and i 3,4
6 | w           : -> e 1, 5
7 | ~w          : and e2 2
8 | ⊥           : neg e 6,7
9 ~r            : neg i 3-8
\end{verbatim}

\paragraph{b)}

\cref{mytree} shows the parse tree of
$(((\neg p) \land q) \to (p \land (q \lor ( \neg r))))$, as shown in
Figure 1.3 on page 34 of [HR04].


\begin{figure}[ht]
  \centering
  \begin{forest}
    for tree={%
      draw,
      circle,
      l sep=1mm,
      s sep=10mm
    }
    [$\to$
      % (\neg p) \land q
      [$\land$
        [$\neg$ [$p$] ]
        [$q$]
      ]
      % p \land (q \lor ( \neg r))
      [$\land$
        [$p$]
        [$\lor$
          [$q$]
          [$\neg$ [$r$] ]
        ]
      ]
    ]
  \end{forest}
  \caption{A parse tree representing a well-formed formula.}
  \label{mytree}
\end{figure}

\paragraph{c)}

\cref{mytable} shows the truth table of $(p \to \neg q) \to (q \lor \neg p)$
and all of its subformulas, as shown in Figure 1.8 on page 40 of [HR04].
The semantics of formulas can be written as $\sem{\varphi}_{v}$.
Semantical entailment is written $\Gamma \vDash \varphi$.
Models for first-order logic are written as $\Model$.

Please always use the \textbf{most-significant-bit-first order} to produce truth-table,
as it is done in \cref{mytable}.

\begin{table}[ht]
    \centering
    \begin{tabular}[t]{cc|cccc|c}
        \toprule
        \rotatebox{90}{$p$}
        & \rotatebox{90}{$q$}
        & \rotatebox{90}{$\neg p$}
        & \rotatebox{90}{$\neg q$}
        & \rotatebox{90}{$p \to \neg q$}
        & \rotatebox{90}{$q \lor \neg p$}
        & \rotatebox{90}{$(p \to \neg q) \to (q \lor \neg p)$} \\
        \midrule
        \FF & \FF & \TT  & \TT  & \TT        & \TT         & \TT \\
        \FF & \TT & \TT  & \FF  & \FF        & \TT         & \TT \\
        \TT & \FF & \FF  & \TT  & \TT        & \FF         & \FF \\
        \TT & \TT & \FF  & \FF  & \TT        & \TT         & \TT \\
        \bottomrule
    \end{tabular}
    \caption{An example of a truth table for a more complex logical formula.}
    \label{mytable}
  \end{table}

\section*{Question 2}
\paragraph{a)}

A table of logical connectives and atoms:
\begin{center}
  \begin{tabular}{l|c}
    Name & Connective \\\hline
    Conjunction (and) & $\phi \land \psi$ \\
    Disjunction (or) & $\phi \lor \psi $ \\
    Implication & $\phi \to \psi $ \\
    Logical equivalence & $\phi \bimpl \psi $ \\
    Negation & $\neg \phi$ \\
    Absurdity (bottom) & $\bot$ \\
    Truth (top) & $\top$ \\
    Universal quantification & $\all{x} \phi$ \\
    Existential quantification & $\exist{x} \phi$ \\
    Uniqueness quantification & $\exist1{x} \phi$ \\
    Equality & $\lequal$ \\
    Inequality & $\lnequal$
  \end{tabular}
\end{center}


\paragraph{b)}
The following is shows a derivation of $\seq{p \land q \to r}{p \to q \to r}$,
where $\Gamma = p \land q \to r, p, q$ in \textbf{proof tree style}.
\begin{equation*}
  \begin{prooftree}
    \infer0{}
    \infer1[\Ax]{\Gamma \vdash p \land q \to r}
    \infer0{}
    \infer1[\Ax]{\Gamma \vdash q}
    \infer0{}
    \infer1[\Ax]{\Gamma \vdash p}
    \infer2[\IAnd]{\Gamma \vdash p \land q}
    \infer2[\EImpl]{\Gamma \vdash r}
    \infer1[\IImpl]{p \land q \to r, p \vdash q \to r}
    \infer1[\IImpl]{p \land q \to r \vdash p \to q \to r}
  \end{prooftree}
\end{equation*}

\paragraph{c)}
The following show a derivation of $p \lor q, \neg q \vdash p$ in
Fitch-style \ND{}.
The full package and documentation can by found here:
\begin{center}
\texttt{https://www.mathstat.dal.ca/\textasciitilde selinger/fitch/}.
\end{center}
Note that this template comes with an adoption of Selinger's package to use
the notation of the lecture.
Further note that the command \texttt{\textbackslash{}r} (reiteration) in the
documentation corresponds to the command \texttt{\textbackslash{}assum} of
the system \ND{} here.

\begin{equation*}
  \begin{nd}
    \hypo {1} {p \lor q}
    \hypo {2} {\neg q}
    \open
    \hypo {3} {p}
    \have {4} {p}        \assum{3}
    \close
    \open
    \hypo {aa} {q}
    \have {6} {\neg q}   \assum{2}
    \have {7} {\bot}     \ne{aa,6}
    \have {8} {p}        \be{7}
    \close
    \have {9} {p}        \oe{1,3-4,aa-8}
  \end{nd}
\end{equation*}

The following shows a proof of $\seq{\neg\neg\phi}{\phi}$ by contradiction in
system \cND{}.
\begin{equation*}
  \begin{nd}
    \hypo {1} {\neg\neg\phi}
    \open
    \hypo {2} {\neg \phi}
    \have {3} {\bot}         \ne{1,2}
    \close
    \have{4} {\phi}          \contra{2-3}
  \end{nd}
\end{equation*}

The following shows a proof of $\seq{\neg \exist{x} P(x)}{\all{y} \neg P(y)}$.
Note that we have to use \texttt{\textbackslash have[][2]\{\}\{\}} to typeset $P(y)$ on line 3 rather than line 2.
This is necessary to show that the introduction of the variable $y$ into context and the assumption of $P(y)$
each start new subproofs.
Also note that the style by Selinger does not explicitly annotate \EAll{} and \IEx{} with substitutions and
we leave them thus out.
\begin{equation*}
  \begin{nd}
    \hypo {1} {\neg \exist{x} P(x)}
    \open[y]
    \have[][2]{}{}
    \open
    \hypo {3} {P(y)}
    \have {4} {\exist{y} P(y)}       \Ei{3}
    \have {5} {\bot}                 \ne{1,4}
    \close
    \have {6} {\neg P(y)}            \ni{3-5}
    \close
    \have {7} {\all{y} \neg P(y)}    \Ai{3-6}
  \end{nd}
\end{equation*}

The following shows a proof of $\seq{}{\all{x} x \lequal x}$ in \NDFe{}, using the reflexivity rule.
The rules \Repl{}, \Sym{} and \Trans{} can be used in Fitch-style proofs using
\texttt{\textbackslash{}repl}, \texttt{\textbackslash{}symm}, and \texttt{\textbackslash{}trans}.
Note the small first letter and that symm has two ``m''!

\begin{equation*}
  \begin{nd}
    \open[x]
    \have[][1]{}{}
    \have {2} {x \lequal x}         \refl{}
    \close
    \have {3} {\all{x} x \lequal x} \Ai{1-2}
  \end{nd}
\end{equation*}

For \textbf{uniform proofs} you can use \texttt{\textbackslash{}backch} for the backchaining rule
\begin{equation*}
  \begin{nd}
    \hypo {h} {\Gamma}
    \have {s1} {\dotsm}
    \have {s2} {\dotsm}
    \have {g} {P(s, t)} \backch{h,s1,s2}
  \end{nd}
\end{equation*}
and $\ufseq[\Gamma]{\phi}$ for sequents that are provable by a uniform proof.
\end{document}